\section{Inverse problem: Anomaly detection from boundary measurements}

\subsection{Framework and numerical approach}

In this section, we study an anomaly corresponding to a single unknown high-intensity heat source $f = \alpha \mathbf{1}_{C}$. Both its intensity $\alpha$ and its physical location $(x_c, y_c)$ (the center of the circular region $C$ of radius $r_0$) are unknown parameters that must be identified. 

To perform this detection, we rely on temperature measurements $T^{\text{mes}}$ taken on the lateral boundaries of the oven $\Gamma_{\text{side}}$. 

For a fixed position $(x_c, y_c)$, we seek the intensity $\alpha$ that minimizes the $L^2$-error between our model prediction and the measurements on $\Gamma_{\text{side}}$:
\begin{equation} \label{eq:cost_inverse}
\mathcal{J}(x_c, y_c, \alpha) = \frac{1}{2}\int_{\Gamma_{\text{side}}} \left|T(x_c,y_c,\alpha) - T^{\text{mes}}\right|^2\dd\sigma.
\end{equation}

Minimizing this functional directly with respect to both the spatial coordinates and the intensity is a difficult problem. To overcome this difficulty, we apply the exact same strategy as in \cref{sec:opti_pb}. Using the linearity of the stationary heat equation, the temperature field can be decomposed using the superposition principle:
\begin{equation}
    T(x_c, y_c, \alpha) = T_0 + \alpha\,\tilde{T}(x_c, y_c),
\end{equation}
where $T_0$ is the solution to~\cref{eq:T0} (without any anomaly) and $\tilde{T}$ follows~\cref{eq:Ti} (with a unit source ($\alpha=1$) strictly located at $(x_c, y_c)$).

The necessary optimality condition $\frac{\partial \mathcal{J}}{\partial \alpha} = 0$ provides an explicit analytical formula for the optimal intensity associated with any given position:
\begin{equation} \label{eq:alpha_opt_inverse}
\alpha^*(x_c,y_c) = \frac{\displaystyle\int_{\Gamma_{\text{side}}} (T^{\text{mes}}-T_0)\,\tilde{T}\dd\sigma}{\displaystyle\int_{\Gamma_{\text{side}}} \tilde{T}^2\dd\sigma}.
\end{equation}

By injecting $\alpha^*(x_c,y_c)$ back into \cref{eq:cost_inverse}, the cost functional depends solely on the spatial coordinates.

We proceed with a grid search over a $30\times23$ spatial grid covering the inner viable domain $[-1.45, 1.45]\times[-0.95, 0.95]$. For each grid point $(x_k, y_k)$:
\begin{enumerate}
    \item We solve the direct problem to obtain $\tilde{T}(x_k, y_k)$ using~\cref{eq:Ti}.
    \item We compute the optimal intensity $\alpha^*(x_k, y_k)$ using \cref{eq:alpha_opt_inverse}.
    \item We evaluate the reduced cost functional $\mathcal{J}(x_k, y_k, \alpha^*)$.
\end{enumerate}
The grid point that minimizes $\mathcal{J}$ is identified as the anomaly's location, and its associated $\alpha^*$ is the detected intensity.

We implemented this approach in \texttt{inversePb.edp}.
For different values of the true anomaly parameters $(x_c, y_c, \alpha)$, we successfully detected the anomaly with good accuracy. A slight discrepancy remains, which is inherently tied to the spatial discretization limit of the grid. While a nested two-pass strategy (performing a localized second grid search around the initially detected position) could further reduce this error, the single-pass high-resolution grid already yields highly satisfactory accuracy while maintaining a reasonable computational time.
